\subsection*{Аналитическое решение}

Казалось бы, у нас есть формула и чётко поставленная задача "--- почему бы нам не решить это \textit{аналитически?} (то есть в явном виде вывести формулу).

Рассмотрим задачу оптимизации для \texttt{MSE}:
\[
\frac{1}{\ell} \sum_{i=1}^\ell (y_i - a(x_i))^2 \rightarrow min_{w}
\]

т.е. необходимо подобрать такие параметры $w$, при которых значение выражения слева будет минимально возможным.

Далее будем работать с матричной формой этого выражения. Напомним, что в матричной форме:
\begin{enumerate}
    \item $X$ "--- матрица объекты-признаки;
    \item $w$ "--- вектор весов;
    \item $y$ "--- вектор реальных ответов.
\end{enumerate}

Итак, перепишем это дело в матричной форме:
\[
\frac{1}{\ell} (y - Xw)^T (y - Xw)
\]

Для дальнейших рассуждений, опустим коэффициент $\frac{1}{\ell}$
и раскроем скобки:
\[
y^Ty - y^TXw - X^Tw^Ty + X^Tw^TXw
\]

Второй и третий члены можем объединить в 1:
\[
y^Ty - 2w^TX^Ty + X^Tw^TXw
\]

Логично предположить, что раз мы ищем минимум функции "--- нужно взять производную и приравнять её к нулю. Однако здесь перед нами встаёт следующая проблема: ранее нам приходилось брать производную по каким"=то неизвестным одномерным величинам, однако сейчас это векторы.

Тут стоит упомянуть довольно сложное т.н. дифференциирование по векторам. В рамках данной лекции мы не будем нагружать излишней теорией, а лишь скажем, что производная выражения выше будет следующей:
\[
\frac{ \partial L }{ \partial w } = -2X^Ty + 2X^TXw
\]
Выражение слева означает \textit{частную производную}. В нашей формуле применяется сразу 3 буквенных обозначения, и тут частной производной мы задаём, по какому аргументу брать производную. В нашем случае этот аргумент "--- вектор весов. В таком случае, в выражении справа всё, что не $w$, считается константой.

Теперь, приравниваем к $0$:
\[
-2X^Ty + 2X^TXw = 0
\]
\[
2X^TXw = 2X^Ty
\]
\[
w = (X^TX)^{-1}X^Ty
\]

Ну и соответственно таким образом мы вывели \textit{аналитическое решение} для \texttt{MSE}. Подставив и вычислив, мы получим наилучший вектор весов для нашей задачи. 

Казалось бы "--- подставляй и считай, однако и здесь не без проблем:
\begin{enumerate}
    \item Операция обращения матрицы свидетельствует о том, что матрица $X$ должна быть полного ранга, что в каком"=то смысле является доп. ограничением на наши данные;
    \item Не уходя далеко от обращения, в среднем данная операция вычисляется за $O(d^3)$, что на больших данных непозволительно и неэффективно;
    \item Представленная формула "--- оптимальные веса с точки зрения только \texttt{MSE}. Далеко не факт, что удастся получить аналитическое решение для иной, более сложной функции.
\end{enumerate}

Все эти пункты (особенно 2 и 3) наталкивают на мысль о том, что нужно использовать некоторый алгоритм, который позволил бы это делать значительно проще. Поскольку мы только что вывели формулу в явном виде, мы уже получили некоторое знание "--- гарантированно можно решить эту задачу не перебором.

Копая глубже, мы натыкаемся на изящный алгоритм, который называется 
\textit{градиентный спуск}.